The increasing complexity of distributed intelligent systems has created a need for theoretical frameworks capable of describing emergent behavior across large-scale adaptive networks. This thesis introduces the concept of Quantized Echo Dynamics (QED), a novel mathematical framework for modeling information propagation within hypermodular cognitive meshes. Unlike conventional network models that treat node interactions as continuous and memoryless processes, the proposed approach represents communication events as discrete echo states that accumulate and transform across hierarchical modular structures.

A formal representation of hypermodular cognitive meshes is developed using graph-theoretic and stochastic methods. The resulting model captures both local adaptation and global self-organization through a system of quantized feedback operators. Analytical results demonstrate that echo-state quantization improves network resilience by reducing propagation instability while preserving adaptive responsiveness. Furthermore, the framework establishes conditions under which emergent synchronization and distributed consensus can arise without centralized coordination.

To evaluate the proposed theory, a series of computational experiments were conducted using synthetic cognitive networks ranging from 10$^3$ to 10$^6$ nodes. Simulation results indicate that QED-based architectures exhibit significantly greater fault tolerance and recovery capability than traditional diffusion and message-passing approaches. In particular, the model maintained stable information coherence under high levels of node failure and communication noise.

The findings suggest that quantized echo dynamics provide a promising foundation for the design of next-generation autonomous infrastructures and large-scale intelligent systems. Beyond their immediate application to distributed computing, the principles developed in this work may contribute to future research in adaptive networks, collective intelligence, and complex systems theory.